\documentclass[12pt]{article}
\usepackage{amsmath,amssymb,amsfonts}
\usepackage[margin=1in]{geometry}
\usepackage{bm}

\newcommand{\vect}[1]{\mathbf{#1}}
\newcommand{\grad}{\boldsymbol{\nabla}}

\begin{document}

\section*{Problem 6}

\textbf{Problem:} Prove the vector identities:
\begin{itemize}
\item[(a)] $\grad(\vect{A}\cdot\vect{B}) = (\vect{A}\cdot\grad)\vect{B} + (\vect{B}\cdot\grad)\vect{A} + \vect{A}\times(\grad\times\vect{B}) + \vect{B}\times(\grad\times\vect{A})$
\item[(b)] $\grad\cdot(\vect{U}\times\vect{V}) = \vect{V}\cdot(\grad\times\vect{U}) - \vect{U}\cdot(\grad\times\vect{V})$
\end{itemize}

\bigskip
\textbf{Solution:}

\subsection*{(a) $\grad(\vect{A}\cdot\vect{B})$}

We work in Cartesian coordinates using index notation with the Einstein summation convention. The $i$-th component of $\grad(\vect{A}\cdot\vect{B})$ is:
\[
[\grad(\vect{A}\cdot\vect{B})]_i = \partial_i(A_j B_j) = (\partial_i A_j)B_j + A_j(\partial_i B_j).
\]

We need to show this equals the $i$-th component of the right-hand side. Consider each term:

\textbf{Term 1:} $[(\vect{B}\cdot\grad)\vect{A}]_i = B_j \,\partial_j A_i$.

\textbf{Term 2:} $[(\vect{A}\cdot\grad)\vect{B}]_i = A_j \,\partial_j B_i$.

\textbf{Term 3:} $[\vect{B}\times(\grad\times\vect{A})]_i = \varepsilon_{ijk}\,B_j\,(\grad\times\vect{A})_k = \varepsilon_{ijk}\,B_j\,\varepsilon_{k\ell m}\,\partial_\ell A_m$.

Using the identity $\varepsilon_{ijk}\varepsilon_{k\ell m} = \delta_{i\ell}\delta_{jm} - \delta_{im}\delta_{j\ell}$:
\[
[\vect{B}\times(\grad\times\vect{A})]_i = (\delta_{i\ell}\delta_{jm} - \delta_{im}\delta_{j\ell})\,B_j\,\partial_\ell A_m
= B_j\,\partial_i A_j - B_j\,\partial_j A_i.
\]

Similarly, \textbf{Term 4:}
\[
[\vect{A}\times(\grad\times\vect{B})]_i = A_j\,\partial_i B_j - A_j\,\partial_j B_i.
\]

Adding all four terms:
\begin{align}
&[(\vect{B}\cdot\grad)\vect{A}]_i + [(\vect{A}\cdot\grad)\vect{B}]_i + [\vect{B}\times(\grad\times\vect{A})]_i + [\vect{A}\times(\grad\times\vect{B})]_i \\
&= B_j\partial_j A_i + A_j\partial_j B_i + (B_j\partial_i A_j - B_j\partial_j A_i) + (A_j\partial_i B_j - A_j\partial_j B_i) \\
&= B_j\partial_i A_j + A_j\partial_i B_j \\
&= \partial_i(A_j B_j) = [\grad(\vect{A}\cdot\vect{B})]_i. \quad \blacksquare
\end{align}

\subsection*{(b) $\grad\cdot(\vect{U}\times\vect{V})$}

In index notation:
\[
\grad\cdot(\vect{U}\times\vect{V}) = \partial_i(\vect{U}\times\vect{V})_i = \partial_i(\varepsilon_{ijk}\,U_j V_k) = \varepsilon_{ijk}\bigl[(\partial_i U_j)V_k + U_j(\partial_i V_k)\bigr].
\]

Consider the first term:
\[
\varepsilon_{ijk}(\partial_i U_j)V_k = V_k\,\varepsilon_{kij}\,\partial_i U_j = V_k\,(\grad\times\vect{U})_k = \vect{V}\cdot(\grad\times\vect{U}),
\]
where we used the cyclic property $\varepsilon_{ijk} = \varepsilon_{kij}$ and recognized the $k$-th component of the curl.

For the second term:
\[
\varepsilon_{ijk}\,U_j(\partial_i V_k) = U_j\,\varepsilon_{jki}\,\partial_i V_k = -U_j\,\varepsilon_{jik}\,\partial_i V_k = -U_j\,(\grad\times\vect{V})_j = -\vect{U}\cdot(\grad\times\vect{V}),
\]
where we used $\varepsilon_{ijk} = \varepsilon_{jki}$ (cyclic) and $\varepsilon_{jki} = -\varepsilon_{jik}$ (swap of last two indices).

Therefore:
\[
\boxed{\grad\cdot(\vect{U}\times\vect{V}) = \vect{V}\cdot(\grad\times\vect{U}) - \vect{U}\cdot(\grad\times\vect{V})}. \quad \blacksquare
\]

\end{document}
